\documentclass[conference]{IEEEtran}
\IEEEoverridecommandlockouts
% The preceding line is only needed to identify funding in the first footnote. If that is unneeded, please comment it out.
%Template version as of 6/27/2024

\usepackage{cite}
\usepackage{amsmath,amssymb,amsfonts}
\usepackage{algorithmic}
\usepackage{graphicx}
\usepackage{textcomp}
\usepackage{xcolor}
\usepackage[utf8]{vietnam}
\usepackage[vietnamese, english]{babel}
\usepackage{booktabs}
\usepackage{tabularx}
\usepackage{array}
\usepackage{multirow}
\usepackage{url}
\def\BibTeX{{\rm B\kern-.05em{\sc i\kern-.025em b}\kern-.08em
    T\kern-.1667em\lower.7ex\hbox{E}\kern-.125emX}}

% Đánh dấu phần chưa viết bằng chữ đỏ trong PDF, để không sót khi nộp bài
\newcommand{\todo}[1]{\textcolor{red}{[TODO: #1]}}

\begin{document}

\title{Báo cáo bài tập lớn: Học máy với PyTorch theo CS188 Project 5 (UC Berkeley)}

\author{\IEEEauthorblockN{Lương Thế Vinh}
\IEEEauthorblockA{\textit{Mã sinh viên:} 23020715}
\and
\IEEEauthorblockN{Nguyễn Trọng Khang}
\IEEEauthorblockA{\textit{Mã sinh viên:} 24020169}
\and
\IEEEauthorblockN{Nguyễn Minh Đức}
\IEEEauthorblockA{\textit{Mã sinh viên:} 24022628}
}

\maketitle

\renewcommand{\abstractname}{Tóm tắt}
\renewcommand{\figurename}{Hình}
\renewcommand{\tablename}{Bảng}
\renewcommand{\IEEEkeywordsname}{Từ khoá}
\renewcommand{\refname}{Tài liệu tham khảo}

\begin{abstract}
Báo cáo trình bày bài tập lớn môn Trí tuệ nhân tạo (INT3401\_11) của nhóm,
thực hiện Project 5: Machine Learning của học phần CS188, Đại học California,
Berkeley. Nhóm tự cài đặt và huấn luyện bằng PyTorch sáu mô hình có độ phức
tạp tăng dần: perceptron phân loại nhị phân, mạng truyền thẳng hồi quy hàm sin,
mạng nhiều lớp phân loại chữ số viết tay MNIST, mạng nơ-ron hồi quy nhận diện
ngôn ngữ của một từ, lớp tích chập tự cài đặt cho bài toán MNIST, và khối
self-attention có mặt nạ nhân quả. Mỗi mô hình được chấm bằng autograder của
dự án theo một ngưỡng độ chính xác hoặc giá trị mất mát.
\todo{một câu tóm tắt kết quả chính, viết sau khi xong phần IV}.
Mã nguồn và hướng dẫn chạy lại được công bố tại
\url{https://github.com/lgthevinh/uet-int3401}.
\end{abstract}

\begin{IEEEkeywords}
perceptron, mạng nơ-ron, MNIST, mạng nơ-ron hồi quy, tích chập, self-attention, PyTorch
\end{IEEEkeywords}

\begin{table}[!ht]
    \centering
    \renewcommand{\arraystretch}{1.15}
    \footnotesize
    \caption{Bảng phân công công việc cho từng thành viên.}
    \label{tab:assignment}
    \begin{tabularx}{\columnwidth}{@{}lccX@{}}
    \toprule
        \textbf{Thành viên} & \textbf{MSSV} & \textbf{Đóng góp} & \textbf{Công việc} \\
    \midrule
        Lương Thế Vinh & 23020715 & 36\% & Nhóm trưởng; Q3, Q5; tổng hợp báo cáo; quản lý tiến độ trên CoReach \\
        Nguyễn Trọng Khang & 24020169 & 32\% & Q1, Q2 \\
        Nguyễn Minh Đức & 24022628 & 32\% & Q4, Q6 \\
    \bottomrule
    \end{tabularx}
\end{table}

% ============================================================
\section{Giới thiệu}
\label{sec:intro}

Bài tập lớn của nhóm thực hiện Project~5: Machine Learning của học phần CS188
(Introduction to Artificial Intelligence) tại Đại học California, Berkeley,
phiên bản PyTorch học kỳ Fall 2025~\cite{cs188}. Dự án gồm sáu bài (Q1--Q6).
Ở mỗi bài, nhóm tự cài đặt mô hình, hàm mất mát và vòng huấn luyện trong các
tệp \texttt{models.py}, \texttt{losses.py} và \texttt{train.py}, chỉ dùng các
thành phần cơ bản của PyTorch~\cite{pytorch}. Kết quả được chấm bằng autograder
đi kèm dự án theo một ngưỡng cố định cho từng bài (Bảng~\ref{tab:overview}),
với tổng điểm tối đa là 25.

Mục tiêu của dự án là hiểu cách một mạng nơ-ron được xây dựng và huấn luyện từ
các phép toán tensor, thay vì dùng các mô hình có sẵn. Sáu bài được sắp theo độ
phức tạp tăng dần, mỗi bài thêm một ý tưởng so với bài trước:
\begin{itemize}
  \item Q1: mô hình tuyến tính đơn giản nhất, perceptron~\cite{rosenblatt},
  phân loại nhị phân bằng một siêu phẳng.
  \item Q2: thêm hàm kích hoạt phi tuyến để xấp xỉ một hàm liên tục
  ($\sin x$) bằng mạng truyền thẳng.
  \item Q3: chuyển từ hồi quy sang phân loại đa lớp, áp dụng cho ảnh chữ số
  viết tay MNIST~\cite{lecun}.
  \item Q4: xử lý đầu vào có độ dài thay đổi (chuỗi ký tự) bằng mạng nơ-ron
  hồi quy (RNN)~\cite{elman}.
  \item Q5: khai thác cấu trúc không gian của ảnh bằng phép tích chập tự cài
  đặt~\cite{lecun}.
  \item Q6: cài đặt khối self-attention có mặt nạ nhân quả, thành phần cốt lõi
  của kiến trúc Transformer~\cite{vaswani}.
\end{itemize}

\begin{table}[!t]
    \centering
    \renewcommand{\arraystretch}{1.15}
    \footnotesize
    \caption{Tổng quan các bài toán và điều kiện đạt của autograder.}
    \label{tab:overview}
    \begin{tabularx}{\columnwidth}{@{}cXXc@{}}
    \toprule
        \textbf{Bài} & \textbf{Nhiệm vụ} & \textbf{Điều kiện đạt} & \textbf{Điểm} \\
    \midrule
        Q1 & Phân loại nhị phân bằng perceptron
           & Độ chính xác 100\% trên tập huấn luyện & 5 \\
        Q2 & Hồi quy $\sin x$ trên $[-2\pi, 2\pi]$
           & Loss $\le 0.02$ trên tập huấn luyện & 5 \\
        Q3 & Phân loại chữ số MNIST (10 lớp)
           & Độ chính xác $\ge 97\%$ trên tập kiểm tra & 5 \\
        Q4 & Nhận diện ngôn ngữ của một từ (5 ngôn ngữ) bằng RNN
           & Độ chính xác $\ge 81\%$ trên tập kiểm tra & 5 \\
        Q5 & Tích chập tự cài đặt và phân loại MNIST
           & \texttt{Convolve} khớp phép tích chập tham chiếu; độ chính xác
             $\ge 80\%$ trên tập kiểm tra & 3 \\
        Q6 & Self-attention có mặt nạ nhân quả
           & Đầu ra khớp đáp án tham chiếu & 2 \\
    \midrule
           & \textbf{Tổng} & & \textbf{25} \\
    \bottomrule
    \end{tabularx}
\end{table}

Phần còn lại của báo cáo được tổ chức như sau. Phần~\ref{sec:background} tóm
tắt kiến thức nền tảng dùng chung cho các bài. Phần~\ref{sec:method} trình bày
phương pháp và cách cài đặt của từng bài. Phần~\ref{sec:results} mô tả thiết
lập thực nghiệm, kết quả và phân tích. Phần~\ref{sec:conclusion} kết luận và
nêu hướng phát triển.

% ============================================================
\section{Kiến thức nền tảng}
\label{sec:background}

Phần này tóm tắt các khái niệm dùng chung cho nhiều bài, để
phần~\ref{sec:method} chỉ cần trình bày phần riêng của từng bài. Trình bày chi
tiết hơn có thể tham khảo trong~\cite{goodfellow}.

\subsection{Tầng tuyến tính và mạng nhiều lớp}
Một tầng tuyến tính biến đổi đầu vào $x \in \mathbb{R}^{n}$ thành
$u = Wx + b$, với ma trận trọng số $W \in \mathbb{R}^{m \times n}$ và vector
bias $b \in \mathbb{R}^{m}$ là các tham số học được. Mạng nhiều lớp
(multilayer perceptron, MLP) xếp chồng các tầng tuyến tính, xen giữa là một
hàm kích hoạt phi tuyến $\sigma$:
\begin{equation}
h = \sigma(W_1 x + b_1), \qquad \hat{y} = W_2 h + b_2.
\label{eq:mlp}
\end{equation}
Hàm phi tuyến là bắt buộc: nếu bỏ $\sigma$, hai tầng liên tiếp
$W_2(W_1 x + b_1) + b_2$ vẫn chỉ là một ánh xạ tuyến tính, nên mạng không mạnh
hơn một tầng duy nhất. Ngược lại, chỉ với một lớp ẩn đủ rộng, MLP đã có thể
xấp xỉ một hàm liên tục bất kỳ trên một miền bị chặn với sai số tuỳ
ý~\cite{hornik}; kết quả này đúng với mọi hàm kích hoạt không phải đa thức,
trong đó có ReLU~\cite{leshno}. Đây là cơ sở để Q2 xấp xỉ được hàm $\sin x$.

\subsection{Hàm kích hoạt}
Báo cáo dùng ba hàm kích hoạt:
\begin{itemize}
  \item \emph{Hàm dấu}: $\mathrm{sign}(u) = +1$ nếu $u \ge 0$ và $-1$ nếu
  ngược lại, dùng cho dự đoán của perceptron ở Q1. Đạo hàm của hàm này bằng 0
  ở hầu hết mọi nơi nên không dùng được để lan truyền gradient.
  \item \emph{ReLU}: $\mathrm{ReLU}(u) = \max(0, u)$, dùng ở các lớp ẩn từ Q2
  đến Q5. ReLU tính nhanh và có gradient bằng 1 trên miền dương, nên ít bị
  triệt tiêu gradient hơn sigmoid hay tanh.
  \item \emph{Softmax}: biến vector điểm số (logits) $z \in \mathbb{R}^{K}$
  thành phân phối xác suất trên $K$ lớp,
  \begin{equation}
  p_k = \mathrm{softmax}(z)_k = \frac{e^{z_k}}{\sum_{j=1}^{K} e^{z_j}},
  \label{eq:softmax}
  \end{equation}
  dùng cho các bài phân loại (Q3--Q5) và để tính trọng số attention ở Q6.
\end{itemize}

\subsection{Hàm mất mát}
Hàm mất mát đo độ lệch giữa dự đoán và nhãn thật trên một batch gồm $B$ mẫu.
Bài toán hồi quy (Q2) dùng sai số bình phương trung bình (MSE):
\begin{equation}
L_{\mathrm{MSE}} = \frac{1}{B} \sum_{i=1}^{B} (\hat{y}_i - y_i)^2.
\label{eq:mse}
\end{equation}
Các bài phân loại (Q3--Q5) dùng cross-entropy. Với nhãn one-hot $y_i$ và xác
suất dự đoán $p_i$ tính theo~\eqref{eq:softmax},
\begin{equation}
L_{\mathrm{CE}} = -\frac{1}{B} \sum_{i=1}^{B} \sum_{k=1}^{K} y_{ik} \log p_{ik},
\qquad
\frac{\partial L_{\mathrm{CE}}}{\partial z_{ik}} = \frac{p_{ik} - y_{ik}}{B}.
\label{eq:ce}
\end{equation}
Gradient theo logits có dạng gọn: lớn khi mô hình dự đoán sai một cách tự tin,
và giảm dần về 0 khi lớp đúng đã được tách biệt rõ. Hàm
\texttt{cross\_entropy} của PyTorch nhận trực tiếp logits và tự áp dụng softmax
bên trong, nên các mô hình phân loại trả về logits thay vì xác suất.

\subsection{Tối ưu bằng gradient descent}
Huấn luyện là tìm bộ tham số $\theta$ cực tiểu hoá hàm mất mát $L$. Gradient
descent lặp lại phép cập nhật
\begin{equation}
\theta \leftarrow \theta - \eta \nabla_{\theta} L,
\label{eq:gd}
\end{equation}
với $\eta$ là learning rate. Thay vì tính gradient trên toàn bộ tập dữ liệu,
mini-batch SGD tính gradient trên một batch $B$ mẫu, cho nhiều lần cập nhật hơn
với cùng chi phí tính toán, đổi lại gradient nhiễu hơn. Một lượt duyệt hết tập
huấn luyện gọi là một epoch. Adam~\cite{kingma} điều chỉnh bước cập nhật riêng
cho từng tham số dựa trên trung bình động của gradient và của bình phương
gradient, thường hội tụ ổn định hơn SGD. Nhóm dùng SGD cho Q2, Q3, Q5 và Adam
cho Q4.

Gradient được tính bằng lan truyền ngược (backpropagation)~\cite{rumelhart}:
áp dụng quy tắc chuỗi từ hàm mất mát ngược về từng tham số. PyTorch thực hiện
việc này tự động (autograd): lời gọi \texttt{loss.backward()} duyệt đồ thị
tính toán đã được ghi lại trong lượt forward và tính gradient cho mọi tham số,
sau đó \texttt{optimizer.step()} áp dụng~\eqref{eq:gd}. Vì vậy mọi phép tính
trong mô hình phải là phép toán tensor để giữ được đồ thị này, đây là điều kiện
quan trọng khi tự cài đặt tích chập ở Q5.

Perceptron ở Q1 cũng nằm trong khung này. Với hàm mất mát perceptron
$L(w) = \max(0, -y\, w \cdot x)$, gradient bằng $-y x$ khi mẫu bị phân loại
sai và bằng 0 khi phân loại đúng, nên một bước SGD với learning rate $\eta$
cho đúng quy tắc cập nhật perceptron $w \leftarrow w + \eta\, y x$. Q1 và các
bài sau vì vậy chỉ khác nhau ở mô hình và hàm mất mát.

\subsection{Đánh giá và khả năng tổng quát hoá}
Dữ liệu được chia thành tập huấn luyện (train) để cập nhật tham số, tập kiểm
định (validation) để theo dõi quá trình huấn luyện và chọn mô hình, và tập
kiểm tra (test) để đánh giá cuối cùng. Một mô hình có thể đạt loss rất thấp
trên tập huấn luyện nhưng dự đoán kém trên dữ liệu mới (overfitting), nên quá
trình huấn luyện được theo dõi bằng độ chính xác trên tập kiểm định và dừng
sớm (early stopping) khi đạt ngưỡng. Tập kiểm tra chỉ phản ánh đúng khả năng
tổng quát hoá khi không được dùng trong bất kỳ quyết định nào của quá trình
huấn luyện; điều kiện này được xem lại cho từng bài ở
phần~\ref{sec:results}.

% ============================================================
\section{Phương pháp}
\label{sec:method}

Mỗi bài được trình bày theo cùng một khung: bài toán cụ thể, mô hình và công
thức, rồi cách cài đặt. Các kiến thức chung ở phần~\ref{sec:background} không
được nhắc lại, và kết quả của từng bài được trình bày ở
phần~\ref{sec:results}.

\subsection{Q1: Perceptron}
\todo{viết lại từ code (\texttt{PerceptronModel}, \texttt{train\_perceptron});
nguồn \texttt{q1\_perceptron.tex}; khoảng 0.3 trang}

\subsection{Q2: Hồi quy phi tuyến}
\todo{viết lại từ code: MLP $1 \to 100 \to 1$, SGD, learning rate 0.01, batch
10; nguồn \texttt{q2\_regression.tex}; khoảng 0.4 trang}

\subsection{Q3: Phân loại chữ số MNIST}
\todo{rút gọn từ \texttt{q3\_digit\_classification.tex}, bỏ phần softmax và
cross-entropy đã có ở mục II-C; khoảng 0.5 trang}

\subsection{Q4: Nhận diện ngôn ngữ}
\todo{rút gọn từ \texttt{q4\_language\_id.tex}, bỏ phần Adam và overfitting đã
có ở mục II; khoảng 0.6 trang}

\subsection{Q5: Tích chập}
\todo{rút gọn từ \texttt{q5\_convolution.tex}; khoảng 0.6 trang}

\subsection{Q6: Self-attention có mặt nạ nhân quả}
\todo{rút gọn từ \texttt{q6\_attention.tex}, bỏ đoạn nhắc tới Q7; khoảng 0.6
trang}

% ============================================================
\section{Thực nghiệm và kết quả}
\label{sec:results}

\subsection{Thiết lập thực nghiệm}
\todo{môi trường chạy; dữ liệu và cách chia tập của từng bài; Bảng III: siêu
tham số (optimizer, learning rate, batch size, số epoch tối đa, điều kiện
dừng)}

\subsection{Kết quả tổng hợp}
\todo{Bảng IV: chỉ số, ngưỡng, kết quả đạt được và điểm của từng bài; tổng
điểm sau khi chạy \texttt{python autograder.py}}

\subsection{Phân tích}
\todo{Q3: so sánh cross-entropy và MSE; Q4: xáo thứ tự batch và giữ trọng số
tốt nhất; Q5: quá trình huấn luyện và lưu ý về độ chính xác 100\%; Q6: kiểm
tra tính đúng; so sánh MLP với tích chập (Q3, Q5) và RNN với attention (Q4,
Q6)}

% ============================================================
\section{Kết luận và hướng phát triển}
\label{sec:conclusion}
\todo{tóm tắt kết quả; hạn chế (\texttt{Convolve} dùng vòng lặp và một
kernel, RNN thuần, attention một đầu, cách chia dữ liệu của Q5); hướng phát
triển (vector hoá tích chập, CNN nhiều kernel, LSTM/GRU, multi-head attention,
character GPT)}

% ============================================================
\begin{thebibliography}{00}
\bibitem{cs188} University of California, Berkeley, ``CS 188 Project 5:
Machine Learning,'' Fall 2025. [Online]. Available:
\url{https://inst.eecs.berkeley.edu/~cs188/archive/fa25/projects/proj5/}
\bibitem{pytorch} A. Paszke \textit{et al.}, ``PyTorch: An imperative style,
high-performance deep learning library,'' in \textit{Proc. Adv. Neural Inf.
Process. Syst. (NeurIPS)}, 2019, pp. 8024--8035.
\bibitem{rosenblatt} F. Rosenblatt, ``The perceptron: A probabilistic model
for information storage and organization in the brain,'' \textit{Psychol.
Rev.}, vol. 65, no. 6, pp. 386--408, 1958.
\bibitem{lecun} Y. LeCun, L. Bottou, Y. Bengio, and P. Haffner,
``Gradient-based learning applied to document recognition,'' \textit{Proc.
IEEE}, vol. 86, no. 11, pp. 2278--2324, Nov. 1998.
\bibitem{elman} J. L. Elman, ``Finding structure in time,'' \textit{Cogn.
Sci.}, vol. 14, no. 2, pp. 179--211, 1990.
\bibitem{vaswani} A. Vaswani \textit{et al.}, ``Attention is all you need,''
in \textit{Proc. Adv. Neural Inf. Process. Syst. (NeurIPS)}, 2017,
pp. 5998--6008.
\bibitem{goodfellow} I. Goodfellow, Y. Bengio, and A. Courville, \textit{Deep
Learning}. Cambridge, MA, USA: MIT Press, 2016.
\bibitem{hornik} K. Hornik, M. Stinchcombe, and H. White, ``Multilayer
feedforward networks are universal approximators,'' \textit{Neural Netw.},
vol. 2, no. 5, pp. 359--366, 1989.
\bibitem{leshno} M. Leshno, V. Y. Lin, A. Pinkus, and S. Schocken,
``Multilayer feedforward networks with a nonpolynomial activation function
can approximate any function,'' \textit{Neural Netw.}, vol. 6, no. 6,
pp. 861--867, 1993.
\bibitem{kingma} D. P. Kingma and J. Ba, ``Adam: A method for stochastic
optimization,'' in \textit{Proc. Int. Conf. Learn. Represent. (ICLR)}, 2015.
[Online]. Available: \url{https://arxiv.org/abs/1412.6980}
\bibitem{rumelhart} D. E. Rumelhart, G. E. Hinton, and R. J. Williams,
``Learning representations by back-propagating errors,'' \textit{Nature},
vol. 323, no. 6088, pp. 533--536, 1986.
\end{thebibliography}

\end{document}
